\scanpage{17}
\uline{Proposition} Soient $V$ un anneau de valuation discrète,
de corps $K$, $A_K$ une V.A. définie sur $K$,
$X$ un préschéma projectif sur $V$, et $R_1$ sur $V$,
[\deleted{fibres} plat, fibres simples en codim $\leq 1$],
$u:X_K\to A_K$ une application rationnelle qui
``engendre'' $A_K$. Alors $A_K$ ``se réduit bien''
sur $V$, i.e. il existe un $A$ abélien sur
$V$ tel que $A_K$ soit $K$-isom. à $A\otimes_V K$.

\deleted{On sait que} Soit $C$ une ``courbe
générique'' dans $\unsure{V|X}$, \unsure{universellement|nécessairement|successivement} simple sur
$V$. (On sait que cela existe\ldots). \deleted{Considérons
\unsure{l'hom|l'image}} Il est connu que alors
$C_K\unsure{-|\to}A_K$ est un \uline{morphisme}, et
qui engendre $A_K$ [question géom. sur
un corps alg. clos]. Considérons l'hom
correspondant $\mathrm{Pic}^0(A_K)\to\mathrm{Pic}^0(C_K)$, \deleted{qui
est injectif, \unsure{On|Or} sait que $\mathrm{Pic}^0(\unsure{C_K|A_K})$
``se réduit bien'', \unsure{On}} \inserted{et en \unsure{retransposant|transposant}}
$\mathrm{Pic}^0(C_K)\to A_K$
qui est surjectif [question géom. sur
un corps]. Comme $\mathrm{Pic}^0(C_K)$ se
réduit bien grâce \unsure{à|:} $\mathrm{Pic}^0(C)$, il s'ensuit
\deleted{On} par le th. de Koizumi-Shimura
que \deleted{Pic} $A$ se réduit bien.
% archive sheet number: 16
